\section{Exercises}\label{sec:04-tactics:06-exercises:1}

\textbf{Key points.} A goal is hypotheses above a line, a statement to prove below it; a tactic replaces it with zero or more simpler goals. The core loop is: try the cheapest tactic, read why it failed, find structure to split or induct on, locate the specific lemma that matches. \texttt{induction} generates one case (and hypothesis \texttt{ih}) per constructor, exactly mirroring a recursive function over the same type.

\begin{enumerate}
\def\labelenumi{\arabic{enumi}.}
\tightlist
\item
  Prove \texttt{theorem and\_comm\_tac \{P Q : Prop\} (h : P ∧ Q) : Q ∧ P := by ...} using \href{https://lean-lang.org/doc/reference/latest/Tactic-Proofs/Tactic-Reference/}{\texttt{constructor}}, \texttt{h.left}, \texttt{h.right}.
\item
  Prove \texttt{theorem nat\_mul\_zero (n : Nat) : n * 0 = 0 := by rfl} --- check whether \href{https://lean-lang.org/doc/reference/latest/Tactic-Proofs/Tactic-Reference/}{\texttt{rfl}} alone works, and if not, use \href{https://lean-lang.org/doc/reference/latest/Tactic-Proofs/Tactic-Reference/}{\texttt{induction}}.
\item
  Rewrite the \texttt{modus\_ponens} proof from Chapter 3 in tactic mode.
\end{enumerate}

Solutions: \hyperref[sec:14-appendix-solutions:03-chapter-4:1]{Appendix, Chapter 4}.

With definitions, propositions, and tactics in hand, we pause for a brief rigor check before diving into groups. Continue to
